The Gospel according to Luke was written by a companion of Paul, believed to be Luke the Evangelist, around 80-90 AD (approximately 60 years after the death of Christ). The same author also wrote the Book of Acts, making them a logical two volume set. As with the other synoptic gospels, much material is shared and scholars have proposed that all three (Matthew, Mark, and Luke) share a unique source, {\it Q}. This has not been verified though.

Luke was writing to a primarily Gentile audience, unlike Matthew, and could therefore exclude much of the highlighting of Old Testament scripture and its links to Jesus being the Messiah, in favor of a more direct exposition of Christian theology (for example, the Sermon on the Plain doesn't have any of the references to the Law of Moses that the Sermon on the Mount contains).
